\begin{lemma}[Initial bridge]
Let $F$ be a front on $\mathbb N$ with $\emptyset\notin F$.  If
$m<n$, then
\[
  \mathsf{FrontBridge}\bigl(F,\{m\},\{n\},\{m,n\}\bigr).
\]
\end{lemma}
\begin{proof}
By \noderef{FrontTreeSingleton}, both singletons belong to $T(F)$, so the
first two membership clauses in \noderef{FrontBridge} hold.  The set
$u=\{m,n\}$ is nonempty.  Since $m<n$, an element of $u$ has a smaller
element of $u$ exactly when it is $n$.  Therefore the finite-tail
definition \noderef{FrontBridgeTail} gives
\[
 \mathsf{FrontBridgeTail}(\{m,n\})=\{n\},
\]
which is nonempty as well.

Suppose first that $\{m\}\notin F$.  Since $m<n$,
$\operatorname{insert}(n,\{m\})=\{m,n\}$.  By the second alternative in
the definition of \noderef{InitialSegment}, with cutoff $n$,
\[
 \{k\in\{m,n\}\mid k<n\}=\{m\};
\]
the equality uses $m<n$ and $n\not<n$.  Thus
$\mathsf{InitialSegment}(\{m\},\{m,n\})$, and the two finite sets are
unequal because $n$ belongs to the latter but not the former.  Hence,
by the definition \noderef{ProperInitialSegment},
$\mathsf{ProperInitialSegment}(\{m\},\operatorname{insert}(n,\{m\}))$.
Applying \noderef{FrontTreeOrderedExtension} to the singleton tree
membership, the assumption $\{m\}\notin F$, and this ordered proper
extension shows that $\{m,n\}\in T(F)$.  This proves the second (left
extension) clause of \noderef{FrontBridge}, with witness $n$.

If instead $\{m\}\in F$, the left clause required by
\noderef{FrontBridge} is \(\mathsf{InitialSegment}(\{m\},\{m,n\})\),
which was just proved using the cutoff $n$; it is not an equality case.

For the right clauses, if $\{n\}\notin F$, the calculation of the tail
above gives exactly
$\{n\}=\mathsf{FrontBridgeTail}(\{m,n\})$.  If $\{n\}\in F$, the
required initial-segment relation holds by the equality alternative in
\noderef{InitialSegment}, since the right node and the bridge tail are
both $\{n\}$.  Thus all clauses in the definition
\noderef{FrontBridge} hold, and consequently
\[
\mathsf{FrontBridge}(F,\{m\},\{n\},\{m,n\}).
\]
\end{proof}
